\chapter{算法选择、问题分解与问题表述}
\label{chap:algorithm-selection-problem-decomposition}
\chaptercontents

到目前为止，我们主要讨论了并行编程的实践知识，包括 CUDA 编程特性、GPU 架构、性能优化技术、并行模式和应用案例。本章将讨论转向更抽象的概念。我们首先把并行编程概括为一种思维过程：设计或选择并行算法，并把领域问题分解为定义清晰、彼此协调的并行工作单元，使每个工作单元都能由所选算法高效完成。算法选择和问题分解得当时，程序员可以在并行度、工作效率和资源消耗之间取得合适的折中。

要为有挑战性的领域问题创建成功的计算解决方案，通常需要把扎实的领域知识与并行思维技能结合起来。具备较强并行思维技能的并行程序员，不必局限于解决领域科学家交付的领域问题。相反，程序员可以与领域科学家合作，分析并变换领域问题的结构：哪些部分本质上是串行的，哪些部分适合高性能并行执行，以及把前一类部分移入后一类时涉及哪些领域特定的权衡。本章将帮助读者进一步理解并行编程和并行思维的一般原则。

\section{算法选择}
\label{sec:algorithm-selection}

算法是一套逐步执行的过程，其中每一步都被精确定义，并且可以由计算机执行。算法必须具备三个基本性质：确定性、有效可计算性和有限性。确定性是指每一步都被精确说明，对应执行什么没有歧义。有效可计算性是指每一步都可以由计算机执行。有限性是指算法必须保证终止。

对于一个给定问题，通常可以提出多种算法来解决。有些算法所需的计算步骤较少，即算法复杂度较低；有些算法暴露出更高程度的并行执行；有些算法的适用范围更广；还有些算法具有更好的精度或数值稳定性。遗憾的是，往往没有一种算法能在所有这些方面都优于其他算法。因此，并行程序员通常需要针对给定的硬件系统，选择能够取得最佳折中的算法。

在本书的多个章节中，我们已经看到如何评估同一计算的不同算法之间的权衡。对于第 11 章的前缀和计算，我们比较了两种并行前缀和算法，即 Kogge--Stone 算法和 Brent--Kung 算法。分析表明，Brent--Kung 算法的算法复杂度较低：它用更少的操作完成同一计算，因此工作效率更高。不过，我们也看到，Kogge--Stone 算法暴露出的并行度高于 Brent--Kung 算法，所以能够用更少的迭代完成计算。算法复杂度与算法暴露的并行度之间的这种权衡，是并行程序员经常遇到的经典取舍。最佳算法通常取决于目标并行硬件的特性，也取决于能否用混合方法缓解算法较高的复杂度，例如组合两种并行算法，或者通过线程粗化把并行算法与低复杂度的顺序算法结合起来。

对于第 14 章的排序模式，我们比较了三种并行排序算法：奇偶排序、归并排序和基数排序。奇偶排序的算法复杂度最高，但并行化最为直接。基数排序是非比较排序，因此算法复杂度可以低于归并排序；它也非常适合并行化。然而，基数排序的通用性不足，因为它只能用于某些类型的键。归并排序是比较排序，比基数排序更通用，可以用于任何具有定义良好的比较运算符的键。在选择并行算法时，通用性与并行执行效率之间的权衡是并行程序员遇到的另一种取舍。

对于第 21 章的静电势图计算，我们比较了两种算法，即直接库仑求和（Direct Coulomb Summation，DCS）和截断求和。这两种方法都暴露出充足的并行度，但它们体现了算法复杂度与精度之间的经典权衡。在第 21 章中我们看到，截断求和算法只牺牲少量精度，就能显著提高网格计算的执行效率。该技术所解决的挑战在于，DCS 这类完全精确的方法，其计算量随体积的平方增长。对于大体积系统，这种增长会使计算时间过长，即使是大规模并行设备也难以承受。截断求和方法利用了这样的观察：许多网格计算问题建立在物理定律之上，距离网格点很远的粒子或样本对数值结果的贡献很小，可以用复杂度低得多的隐式方法统一处理。算法复杂度与精度之间的这种权衡并非并行编程所独有，顺序实现也会遇到它；但它往往给并行程序员带来额外挑战。下面继续用截断求和算法说明这些挑战。

在顺序计算中，一个简单的截断算法一次处理一个原子。对于每个原子，算法遍历以原子坐标为中心、落在给定半径内的网格点。由于网格点存放在数组中，可以根据坐标方便地建立索引，因此这是一个直接的过程。然而，这种简单做法不容易搬到并行执行上。原因是以原子为中心的分解具有散布式内存访问行为，效果并不好。因此，我们需要基于以网格为中心的分解来实现截断分箱算法：每个线程计算一个网格点的能量值。算法的关键是先依据输入原子的坐标把它们排序到分箱中。每个分箱对应网格空间中的一个盒子，其中包含坐标落在该盒子内的所有原子。对于网格点，我们把“邻域”定义为一组分箱；按照截断距离阈值，这些分箱包含所有可能对该网格点能量值作出贡献的原子。第 21 章描述了为所有网格点高效管理邻域分箱的方法。在该算法中，每个块遍历自己的邻域分箱；块中的所有线程共同扫描这些邻域分箱中的原子，但每个线程独立判断某个原子是否落在自己负责网格点的截断半径内。

分箱还有一个微妙问题：不同分箱最后可能包含不同数量的原子。有些分箱可能包含很多原子，有些分箱甚至没有原子。如果把分箱的大小（容量）都设为能够容纳一个分箱可能拥有的最大原子数，开销会大得难以承受。一个广为人知的解决方案是把分箱大小设为合理的值，通常远小于单个分箱可能拥有的最大原子数。分箱过程维护一个溢出列表。把原子排序到分箱时，如果原子的归属分箱已满，就把该原子加入溢出列表。设备完成用于计算静电势图的截断求和内核后，还需要处理溢出列表中的原子，以补上缺少的贡献。

\section{问题分解}
\label{sec:problem-decomposition}

选择合适的算法后，要用并行计算解决一个问题，还必须以某种方式表述问题，使它能够分解为可以安全地同时求解的子问题。在这种表述和分解下，程序员编写代码并组织数据，从而并发解决这些子问题。在大型计算问题中寻找并行性在概念上往往很简单，但在实践中可能颇具挑战。关键是确定每个线程或线程集合要执行的工作单元，或者要解决的子问题，使问题中固有的并行性得到充分利用。

并行执行时，分解问题最常见的两种策略是输出中心分解和输入中心分解。图 \ref{fig:22-1} 展示了这两种分解策略。顾名思义，在输出中心分解中，求解每个子问题会产生一个或多个输出元素的结果。基于输出中心分解的实现，会为一个线程或一组线程分配一个或多个子问题，使其产生一个或多个输出元素。每个子问题的求解可能需要所有输入元素（例如第 21 章以网格为中心的 DCS），也可能只需要部分输入元素（例如第 5 章的矩阵乘法，以及第 21 章使用分箱的以网格为中心的截断求和）。

相反，在输入中心分解中，求解每个子问题会推导出一个或多个输入元素对最终输出的贡献。基于输入中心分解的实现，会为一个线程或一组线程分配一个或多个输入元素，使其处理这些输入元素并把它们的贡献应用到输出元素。每个子问题的求解可能更新一个输出元素（例如第 9 章的直方图）、一部分输出元素（例如假设存在的第 21 章以原子为中心的截断求和），或者所有输出元素（例如第 21 章以原子为中心的 DCS）。

\begin{figure}[htbp]
  \centering
  \begin{tikzpicture}[x=0.55cm,y=0.55cm,>=Stealth]
    \begin{scope}[shift={(0,0)}]
      \node[font=\small] at (0.2,5.9) {输入};
      \foreach \x in {1,...,8} {\draw[fill=white,draw=gray!70] (\x-0.5,5.5) rectangle +(1,0.55);}
      \draw[fill=cyan!65,draw=cyan!50!black] (1.1,3.15) circle (0.72);
      \draw[fill=cyan!65,draw=cyan!50!black] (3.6,3.15) circle (0.72);
      \node[font=\scriptsize] at (1.1,3.15) {线程 1};
      \node[font=\scriptsize] at (3.6,3.15) {线程 2};
      \foreach \x in {0.5,1.5,2.5,3.5,4.5,5.5,6.5,7.5} {
        \draw[->,teal!80!black] (\x,5.45) -- (1.1,3.82);
        \draw[->,teal!80!black] (\x,5.45) -- (3.6,3.82);
      }
      \foreach \x in {1,...,7} {\draw[fill=cyan!45,draw=gray!70] (\x-0.5,0.55) rectangle +(1,0.55);}
      \node[font=\small] at (0.15,0.82) {输出};
      \foreach \x in {0.5,1.5,2.5} {\draw[->,teal!80!black] (1.1,2.45) -- (\x,1.12);}
      \foreach \x in {3.5,4.5,5.5,6.5} {\draw[->,teal!80!black] (3.6,2.45) -- (\x,1.12);}
      \node[font=\small] at (3.5,-0.15) {(a) 输出中心};
    \end{scope}
    \begin{scope}[shift={(9.2,0)}]
      \node[font=\small] at (0.2,5.9) {输入};
      \foreach \x in {1,...,8} {\draw[fill=cyan!65,draw=cyan!50!black] (\x-0.5,5.5) rectangle +(1,0.55);}
      \draw[fill=cyan!65,draw=cyan!50!black] (1.1,3.15) circle (0.72);
      \draw[fill=cyan!65,draw=cyan!50!black] (3.6,3.15) circle (0.72);
      \node[font=\scriptsize] at (1.1,3.15) {线程 1};
      \node[font=\scriptsize] at (3.6,3.15) {线程 2};
      \draw[->,teal!80!black] (0.5,5.45) -- (1.1,3.82);
      \draw[->,teal!80!black] (2.5,5.45) -- (3.6,3.82);
      \foreach \x in {1,...,7} {\draw[fill=white,draw=gray!70] (\x-0.5,0.55) rectangle +(1,0.55);}
      \node[font=\small] at (0.15,0.82) {输出};
      \foreach \x in {0.5,1.5,2.5,3.5} {\draw[->,teal!80!black] (1.1,2.45) -- (\x,1.12);}
      \foreach \x in {3.5,4.5,5.5,6.5} {\draw[->,teal!80!black] (3.6,2.45) -- (\x,1.12);}
      \node[font=\small] at (3.5,-0.15) {(b) 输入中心};
    \end{scope}
  \end{tikzpicture}
  \caption{问题分解策略：（a）输出中心分解和（b）输入中心分解。依据原书图 22.1 重绘；箭头表示输入贡献在工作单元与输出元素之间的映射。}
  \label{fig:22-1}
\end{figure}

对于某些问题，输入中心分解和输出中心分解可能产生相同的子问题。例如，在第 14 章的奇偶排序中，每个线程负责读取一对输入元素，并产生相应的一对输出元素。在这个例子中，由于输入元素和输出元素之间具有简单直接的关系，把分解看作输入中心还是输出中心并不重要。

虽然两种分解策略会得到相同的执行结果，但在给定硬件系统上可能表现出非常不同的性能。输出中心分解通常表现为聚集内存访问：每个线程把输入值的影响聚集到一个输出值中。图 \ref{fig:22-1}(a) 示意了聚集访问行为。聚集式访问模式在 CUDA 设备上通常更理想，因为线程可以在私有寄存器中累加结果。此外，多个线程可以共享输入值，并有效利用 GPU 缓存或共享内存来节省全局内存带宽。

另一方面，输入中心分解通常表现为散布内存访问：每个线程把一个输入值的影响散布或分配到输出值中。图 \ref{fig:22-1}(b) 示意了散布行为。散布式访问模式在 CUDA 设备上通常不理想，因为多个线程可能同时更新同一个输出值。输出值必须存放在所有相关线程都能写入的内存中；多个线程同时写入输出值时，必须使用原子操作来防止竞态条件和数据丢失。这些原子操作明显慢于输出中心分解中使用的寄存器访问。

不过，除了聚集与散布访问模式之外，决定某个应用更适合输出中心分解、输入中心分解或其他分解时，还要考虑许多因素，包括分解暴露出的并行度、识别哪些输入数据对哪些输出数据有贡献的难易程度、分解造成的负载均衡情况，以及其他取决于应用的因素。因此，不能简单地断言聚集总是优于散布。

输入中心分解和输出中心分解之间的比较，在第 21 章最为明显，因为该章实现并明确比较了两种策略。具体而言，以原子为中心的分解属于输入中心分解，以网格为中心的分解属于输出中心分解。不过，除第 21 章以外，本书讨论的许多计算也都隐含地作出了问题分解的设计决策。下面重新考察这些计算，突出每个例子所选择的问题分解，以及它相对于其他分解为何更有优势。

图像处理（第 3 章）、矩阵乘法（第 3、5、15 章）、卷积（第 7 章）和模板计算（第 8 章）都采用了输出中心分解。也就是说，这些计算中的线程被分配给输出元素（图像像素、矩阵元素或网格点），并遍历对它们有贡献的输入元素。另一种输入中心分解会把线程分配给输入元素，让每个线程遍历该输入元素会贡献到的输出元素，并使用原子操作更新输出。在这些例子中，输出中心分解相对于输入中心分解的明显优势，是采用聚集访问而不是散布访问，从而避免原子操作。另一方面，其他考虑因素都没有使输入中心分解更有吸引力：输出元素数量足够多，可以暴露高度并行性；很容易确定每个输出元素所需的输入元素；并且计算所有输出元素需要的工作量相同，因此不存在负载不均衡。

直方图计算（第 9 章）采用了输入中心分解。也就是说，每个线程被分配给一个输入元素（或一组输入元素），并根据输入值更新输出分箱。由于多个线程可能更新同一个输出分箱，因此需要原子操作（散布访问模式）。另一种输出中心分解会把线程分配给输出分箱，让每个线程搜索映射到该分箱的输入，并据此更新分箱。这种分解会消除原子操作，因为每个分箱只由一个线程更新（聚集访问模式），但会带来许多问题。第一，输出分箱的数量通常远小于输入值的数量，因此并行度会大幅降低。第二，线程在检查输入值之前无法知道哪些输入值映射到自己的输出分箱，所以每个线程都必须遍历每个输入值，工作效率很低。第三，即使线程有办法快速找出映射到输出分箱的输入值，不同分箱映射到的输入值数量也会不同，导致线程之间负载不均衡。所有这些因素使直方图更适合输入中心分解。

\noindent\textbf{译者注：}底本把本段的扫描章节写作第 12 章；本书扫描实际为第 11 章。这里按本书章节结构作最小修正。

归约（第 10 章）和扫描（第 11 章）模式都包含多次迭代。如果逐次迭代考察问题分解，那么归约和扫描都采用了输出中心分解。每个线程负责更新本次迭代输出数组中的一个元素，通过聚集并相加本次迭代输入数组中的两个元素来完成工作。归约或扫描某次迭代的输入中心分解会把每个线程分配给本次迭代输入数组中的一个元素，再由该线程把该元素的值散布到本次迭代输出数组中需要它的元素。由于多个输入元素会贡献给同一个输出元素，这种分解需要使用原子操作来避免竞态条件，并会严重妨碍性能。

不稳定过滤和稳定过滤模式（第 12 章）都采用了输入中心分解。每个线程被分配给一个输入元素，负责判断该输入元素是否应保留；如果应保留，还要确定它在输出数组中的位置。要为不稳定或稳定过滤建立高效的输出中心分解会相当困难：它需要把线程分配给输出数组中的一个元素，让该线程确定应放在此处的输入元素。但这样做要求线程检查大量其他元素，既效率低下，又会在线程之间产生冗余工作。因此，输入中心分解自然更合适。值得注意的是，稳定过滤的总体分解是输入中心的，但其中包含的扫描操作本身是输出中心的。复杂计算包含采用不同分解策略并行化的多个操作，并不罕见。

归并操作（第 13 章）采用了输出中心分解。每个线程被分配给输出数组中的一组元素，并执行二分搜索，找出应复制到这些输出元素中的相应输入元素。许多其他计算偏好输出中心分解是为了避免原子操作，但这项考虑与归并操作无关：归并中的每个输出元素只从一个输入元素得到贡献，因此输入中心的归并不需要原子操作。归并偏好输出中心而不是输入中心的主要原因，是在线程粗化时实现负载均衡。对于输入中心的归并，可以把一个输入数组划分给各线程，让每个线程在另一个输入数组中查找对应元素，再把它们归并到输出数组中。然而，不同线程最终可能需要归并数量差异很大的输入元素，从而在线程之间产生控制流分歧和负载不均衡。因此，并行归并更适合输出中心分解。

排序（第 14 章）根据所用算法采用不同的问题分解。对于奇偶排序，输入中心和输出中心分解是等价的。在奇偶排序的每次迭代中，每个线程被分配给一对元素，并决定是否交换它们。换句话说，每个线程被分配给一对输入元素，处理后产生一对输出元素，其中没有散布或聚集。对于归并排序，每一层主要由归并操作组成，因而采用前面所述的输出中心分解。对于基数排序，每次迭代采用输入中心分解：每个线程被分配给一个输入元素，并确定它在输出数组中的新位置。回忆一下，基数排序中使用的划分操作是稳定过滤模式的推广，而稳定过滤同样采用输入中心分解。

波前模式（第 16 章）采用了输出中心分解。为了计算动态规划矩阵中的每个（输出）元素，分配给该元素的线程会聚集多个此前计算出的（输入）元素的值。相反，波前模式的输入中心分解会让每个线程把某个此前计算出的（输入）元素的值散布给依赖该值的所有后续（输出）元素。由于动态规划矩阵中的多个元素可能贡献给同一个后续元素，这种分解需要原子操作，因此不如输出中心分解有利。

对于稀疏矩阵计算（第 17 章），SpMV/CSR、SpMV/ELL 和 SpMV/JDS 内核都使用输出中心分解：每个线程被分配给输出向量的一个元素，并遍历相应输入行中的非零元素，聚集它们的贡献。另一方面，SpMV/COO 内核使用输入中心分解：每个线程被分配给输入矩阵中的一个非零元素，并通过原子操作更新输出向量中相应的元素，形成散布访问模式。与输出中心内核相比，输入中心的 SpMV/COO 内核需要原子操作，这是它的缺点；但它还具有提取更多并行度和实现更好负载均衡等优点。SpMV/COO 的输出中心分解需要线程搜索所有输入元素，找出对自己输出向量元素有贡献的那些元素，这会大幅增加计算复杂度。SpMV/CSC 内核也使用输入中心分解：每个线程被分配给输入向量的一个元素，遍历输入矩阵的一列，并使用原子操作把贡献散布到相关的输出向量元素。CSC 格式没有方便的访问方式，让输出中心分解遍历某一行中的全部非零元素；要找出对一个输出元素有贡献的所有输入元素，每个线程都必须扫描所有列，检查其中是否存在属于该行的输入元素。归根结底，这种计算的最佳分解取决于输入数据集。此外，混合 ELL--COO 格式展示了混合输出中心分解与输入中心分解可能有利的例子。

对于图遍历（第 18 章），以顶点为中心的推送式实现和以边为中心的实现采用输入中心分解：线程被分配给顶点或边，并更新相邻顶点的层级。相比之下，以顶点为中心的拉取式实现采用输出中心分解，每个线程只更新分配给自己的顶点的层级。在不同分解之间作选择时，散布与聚集访问模式并不是关键因素，因为层级值的更新具有幂等性，不需要原子操作。不过，所提取的并行度和达到的负载均衡会对哪种分解更有利起重要作用，而最佳分解最终取决于输入数据集。读者可以回顾第 18 章，以了解这一主题的更详细处理。

静电势图计算（第 21 章）偏好输出中心分解，因为聚集访问模式能够避免原子操作。\noindent\textbf{译者注：}底本此处把 gather 与 scatter 在避免原子操作方面的优劣关系写反；第 21 章的实现表明，以网格为中心的输出中心分解采用聚集访问，通过让线程在寄存器中累加而避免原子操作，因而译稿作此最小修正。

该章已经深入处理了两种分解策略之间的区别。计算涉及大量输入原子对大量输出网格点的势能贡献。一个现实的分子系统模型通常至少包含数十万个原子和数百万个能量网格点。计算原子对能量网格点的贡献时，每个原子的静电荷信息会被多次使用。我们看到，静电势图计算可以采用以原子为中心（即输入中心）的分解，也可以采用以网格为中心（即输出中心）的分解。在以原子为中心的分解中，每个线程负责计算一个原子对所有网格点的影响；相反，以网格为中心的分解让每个线程计算所有原子对一个网格点的影响。以网格为中心的输出中心分解更好，因为它使用了更有利的聚集访问模式。两种分解暴露出的并行度都足够。对于截断求和算法，以网格为中心的分解通过分箱克服了把输入数据映射到输出数据的困难。

\section{应用级考量——阿姆达尔定律}
\label{sec:application-level-amdahl}

现实应用通常由多个协同工作的模块组成。图 \ref{fig:22-2} 展示了分子动力学应用主要模块的概览。对于系统中的每个原子，应用需要计算施加在该原子上的各种力，例如振动力、转动力和非键力。每种力都用不同的方法计算。在高层次上，程序员需要决定如何组织这些工作。注意，不同模块的工作量可能差异很大。非键力计算通常涉及许多原子之间的相互作用，所需计算远多于振动力和转动力。因此，这些模块往往作为对力数据结构的独立遍次来实现。

程序员需要决定是否值得为 CUDA 设备实现每个遍次。例如，程序员可能认为振动力和转动力计算的工作量不足，不值得在 GPU 设备上执行。这样得到的 CUDA 程序会启动一个内核，为所有网格点计算非键力场，同时继续在主机上计算网格点的振动力和转动力。更新原子位置和速度的模块也可以在主机上运行：它先合并来自主机的振动力和转动力，以及来自设备的非键力；然后利用合并后的力计算新的原子位置和速度。

最终由 GPU 设备完成的工作份额决定并行化取得的应用级加速比。例如，假设非键力计算占原始顺序执行时间的 95\%，并使用 GPU 获得 100 倍加速；再假设应用的其余部分仍在主机上执行，且没有得到加速。应用级加速比为
\[
\frac{1}{5\%+95\%/100}=\frac{1}{5\%+0.95\%}=\frac{1}{5.95\%}\approx 17\text{ 倍}
\]
如果主机和 CUDA 设备的执行可以重叠，使并行部分的执行时间完全隐藏在主机执行时间中，那么应用级加速比上限为
\[
\frac{1}{5\%}=20\text{ 倍}
\]
这正是第 1 章介绍的阿姆达尔定律：并行计算带来的应用加速比受应用顺序部分限制。
在这个例子中，虽然顺序部分很小，只有 5\%，但它仍把应用级加速比限制在 20 倍；非键力计算本身获得了 100 倍加速，并且完全隐藏在主机执行顺序部分的阴影中。这个例子说明，加速大型应用的一个主要挑战是：不值得在 CUDA 设备上并行执行的小活动，其累计执行时间可能成为最终用户所看到的加速比的限制因素。

阿姆达尔定律经常促使人们进行任务级并行化。虽然其中一些较小的活动不值得进行细粒度的大规模并行执行，但当数据集足够大时，让其中一些活动彼此并行执行可能很有价值。一种办法是使用多核主机并行执行这些任务。另一种办法是尝试同时执行多个小内核，每个内核对应一个任务。虽然每个内核可能无法充分利用 GPU，相比 CPU 执行也只能取得有限加速，但让多个小内核通过多条 CUDA stream 在 GPU 上同时执行，可以彼此重叠执行时间，从整体上取得良好加速。CUDA 设备通过 stream 支持任务级并行；第 23 章将讨论这种机制。

\begin{figure}[htbp]
  \centering
  \begin{tikzpicture}[>=Stealth,node distance=0.8cm and 1.2cm]
    \node[draw,minimum width=3.0cm,minimum height=0.7cm,align=center] (neighbor) {邻域列表};
    \node[draw,minimum width=3.0cm,minimum height=0.7cm,align=center,below=of neighbor] (nonbond) {非键力};
    \node[draw,minimum width=5.8cm,minimum height=0.7cm,align=center,below=1.0cm of nonbond] (update) {更新原子位置和速度};
    \node[draw,minimum width=3.2cm,minimum height=0.7cm,align=center,left=1.8cm of update] (vib) {振动力与转动力};
    \node[draw,minimum width=3.0cm,minimum height=0.7cm,align=center,below=of update] (next) {下一时间步};
    \draw[->] (neighbor) -- (nonbond);
    \draw[->] (nonbond) -- (update);
    \draw[->] (vib) -- (update);
    \draw[->] (update) -- (next);
  \end{tikzpicture}
  \caption{分子动力学应用的主要任务。依据原书图 22.2 重绘。}
  \label{fig:22-2}
\end{figure}

\section{问题表述}
\label{sec:problem-formulation}

问题表述可以说是并行应用开发中影响最为深远的方面 \cite{ch22-wing2006computational}。“最好”的并行化方法不是使用通用数值方法，把领域问题映射为能够用本书所学技能实现的计算步骤。这样做往往会导致很高的计算复杂度和（或）有限的并行度。相反，我们通常需要重新思考数值方法，放宽实现上的某些约束。

第 21 章的截断分箱方法就是一个很好的例子。该方法需要领域知识，在精度与大幅降低算法复杂度之间作出权衡；还需要问题分解和优化技能，才能使用计算机科学家设计的分箱技术实现以网格为中心的分解。这种方法可能带来最大的性能优势以及根本性的新发现和新能力。例如，它可能使我们能够对传统方法认为规模超出能力范围的生化系统进行高保真模拟。这种方法是跨学科的，需要计算机科学洞见和领域洞见，但付出的努力是值得的。

\section{批处理：延迟与吞吐量}
\label{sec:batching-latency-throughput}

到目前为止讨论的并行编程优化技术，都是在问题相同的情况下让并行程序运行得更快。例如，第 11 章中的双缓冲、warp 级 shuffle、寄存器分块和 look-back 等技术，能够让代码在处理相同输入时运行得更快，也就是具有更短的延迟。让并行代码在输入数据相同的情况下运行得更快，通常称为延迟优化。虽然让并行代码在同一输入上运行更快是理想目标，但这不是优化并行程序的唯一动机。一般来说，优化并行程序至少有两种动机。

第一种动机最直观：用更少的时间解决给定问题，即处理给定的输入数据集。这一直是本书关注的重点。例如，一家投资公司可能需要在收市后分析所有投资组合的金融风险。当前并行程序可能需要 10 小时才能完成分析，但投资组合管理团队可能要求在 4 小时内完成，以便更及时地作出决策。当并行程序无法满足新的延迟要求时，可以使用优化技术进一步提高速度，使分析在规定时间内完成。在第 20 章中，FlashAttention 技术通过消除对全局内存的冗余访问，降低了给定输入的 Transformer 延迟。

第二种动机不那么直观：改善解决许多问题所需的总延迟。在上面的投资组合分析例子中，投资公司可能出于统计原因，决定用不同参数运行多个分析实例。然而，当前并行程序可能无法在规定时间内运行所有实例。优化当前并行程序、降低解决多个问题实例的总延迟，会带来额外权衡：每个问题的单独延迟与所有问题实例的总延迟之间的权衡。

在第 20 章中，处理用户查询可以看作推理问题的一个实例。当许多用户发出查询时，就存在许多个推理问题实例。把不同用户的查询批处理后，可以让 Transformer 层一起处理这些查询，提高注意力头的算术强度。例如，批处理可以把 QKV 投影中的向量—矩阵乘法转换为矩阵乘法。虽然一次矩阵乘法的执行时间比一次向量—矩阵乘法更长，但它比批处理中所有向量—矩阵乘法执行时间之和短得多。这是因为矩阵乘法复用原始矩阵元素，消除了各次向量—矩阵乘法之间冗余的全局内存访问，大幅提高了算术强度。

也就是说，单个查询经历了更长的 QKV 投影延迟，但整组查询完成所需的总延迟更短。因此，批处理时 QKV 投影的吞吐量为
\[
\frac{\text{查询总数}}{\text{矩阵乘法延迟}}
\]
不批处理时的吞吐量为
\[
\frac{\text{查询总数}}{\text{所有向量—矩阵乘法的总延迟}}
\]
前者高于后者。
单个问题实例的延迟与一批问题实例的吞吐量之间的这种权衡，是优化并行程序时需要考虑的重要方面。读者还应注意，第 1 章已经讨论过同样的权衡：CPU 设计处理每个单独数据元素时延迟较短，而 GPU 设计处理大量数据元素时吞吐量较高。

关于算法、问题分解、问题表述和优化策略，已有大量覆盖广泛主题的文献。要全面介绍所有可用概念和技术，超出了本书的范围。我们在算法、问题分解、问题表述和优化的各个步骤中，讨论了一组具有广泛适用性的技术。虽然这些技术通过 CUDA 实现来展示，但它们能帮助读者建立并行思维的一般基础。我们相信，人类通过自底向上的方式学习效果最好：先在特定编程模型的上下文中学习概念，获得坚实基础，再把知识推广开来。深入体验 CUDA 实现也能帮助我们积累经验，从而学习并行编程和并行思维的概念，其中有些概念甚至并不专属于 GPU。

\section{小结}
\label{sec:algorithm-selection-summary}

总之，我们讨论了并行编程和并行思维的主要步骤，即问题表述、算法选择、问题分解和优化。给定一个计算问题，程序员通常需要从多种算法中作出选择。有些算法在保持相同数值精度的同时实现不同的权衡；另一些算法则牺牲一定程度的精度，以获得可扩展性强得多的运行时间。对于所选算法，不同的问题分解选择会在并行执行期间造成不同程度的线程干扰、暴露出的并行度、负载不均衡，以及其他性能差异。还可以使用批处理技术，以增加单个问题实例的延迟换取问题实例组更高的吞吐量。并行思维技能使并行程序员能够绕过障碍，找到良好的解决方案。

\begin{chapterbibliography}{1}
  \bibitem{ch22-wing2006computational}
  J. M. Wing, Computational thinking, \emph{Communications of the ACM} 49 (3) (2006) 33--35.
\end{chapterbibliography}
